\documentclass[aspectratio=169, table, xcdraw]{beamer}

% Tema y Colores Institucionales
\usetheme{Madrid}
\usecolortheme{default}
\definecolor{ucbblue}{RGB}{0, 51, 102}

\setbeamercolor{structure}{fg=ucbblue}
\setbeamercolor*{palette primary}{fg=white,bg=ucbblue}
\setbeamercolor*{palette secondary}{fg=white,bg=ucbblue!80}
\setbeamercolor*{palette tertiary}{fg=white,bg=ucbblue!90}
\setbeamercolor{title}{fg=white, bg=ucbblue}
\setbeamercolor{frametitle}{fg=white, bg=ucbblue}
\setbeamercolor{item}{fg=ucbblue}

\usepackage[T1]{fontenc}
\usepackage[utf8]{inputenc}
\usepackage{tgpagella}
\usefonttheme{serif}
\usepackage[spanish, es-noshorthands]{babel}
\renewcommand{\tablename}{Tabla}
\renewcommand{\figurename}{Figura}

\usepackage{amsmath, amssymb, bm}
\usepackage{graphicx}
\usepackage{booktabs}
\usepackage[most]{tcolorbox}
\usepackage{tikz}
\usetikzlibrary{shapes.geometric, arrows.meta, positioning, calc}

\title[Consolidación FK]{Consolidación Preevaluación de Cinemática Directa}
\subtitle{Resolución, Diagnóstico y Verificación}
\author[B. Quiroga Turdera]{M.Sc. Ing. Bernardo Quiroga Turdera}
\institute[UCB Tarija]{Universidad Católica Boliviana ``San Pablo'' — Sede Tarija \\ Robótica (IMT-342)}
\date{Semana 5 - Sesión 2}

\begin{document}

\begin{frame}
    \titlepage
\end{frame}

\begin{frame}{Flujo de la Sesión: Consolidación Activa}
    \begin{center}
        \large \textit{Esta sesión no es una clase expositiva. Es una evaluación formativa de su capacidad para ejecutar la cinemática directa de forma independiente.}
    \end{center}
    
    \vspace{0.4cm}
    \begin{center}
    \begin{tikzpicture}[
        node distance=0.6cm and 0.6cm,
        box/.style={rectangle, draw=ucbblue, thick, fill=ucbblue!5, text width=2.2cm, align=center, rounded corners, minimum height=1cm, font=\footnotesize},
        arrow/.style={-Stealth, thick, ucbblue, line width=1.2pt}
    ]
        \node[box] (i1) {\textbf{Intento 1}\\(Independiente)};
        \node[box, right=of i1] (diag) {\textbf{Diagnóstico}\\(Identificar Errores)};
        \node[box, right=of diag] (corr) {\textbf{Corrección}\\(Focalizada)};
        \node[box, right=of corr] (i2) {\textbf{Intento 2}\\(Transferencia)};
        \node[box, fill=green!10, right=of i2] (read) {\textbf{Readiness Check}\\(Evaluación)};
        
        \draw[arrow] (i1) -- (diag);
        \draw[arrow] (diag) -- (corr);
        \draw[arrow] (corr) -- (i2);
        \draw[arrow] (i2) -- (read);
    \end{tikzpicture}
    \end{center}
\end{frame}

\begin{frame}{Problema Diagnóstico 1: Cinemática Planar 3R}
    \begin{columns}
        \begin{column}{0.4\textwidth}
            Un manipulador planar 3R tiene la siguiente tabla DH estándar[cite: 5]:
            
            \vspace{0.2cm}
            \resizebox{\textwidth}{!}{
            \begin{tabular}{c c c c c}
            \toprule
            $i$ & $\theta_i$ & $d_i$ [m] & $a_i$ [m] & $\alpha_i$ \\ \midrule
            1 & $q_1$ & 0 & 0.30 & 0 \\
            2 & $q_2$ & 0 & 0.20 & 0 \\
            3 & $q_3$ & 0 & 0.10 & 0 \\ \bottomrule
            \end{tabular}
            }
            
            \vspace{0.3cm}
            \textbf{Configuración actual:} \\
            $q_1=30^\circ, \quad q_2=-60^\circ, \quad q_3=90^\circ$[cite: 5].
        \end{column}
        
        \begin{column}{0.6\textwidth}
            \begin{center}
            \begin{tikzpicture}[scale=6, thick]
                \draw[->, gray] (0,0) -- (0.6,0) node[right]{$X_0$};
                \draw[->, gray] (0,0) -- (0,0.4) node[above]{$Y_0$};
                
                \coordinate (O) at (0,0);
                \coordinate (J1) at (30:0.3);
                \coordinate (J2) at ($(J1) + (-30:0.2)$);
                \coordinate (J3) at ($(J2) + (60:0.1)$);
                
                \draw[ucbblue, line width=2pt] (O) -- (J1) node[midway, above left]{\small $a_1$};
                \draw[ucbblue!70, line width=2pt] (J1) -- (J2) node[midway, below]{\small $a_2$};
                \draw[ucbblue!40, line width=2pt] (J2) -- (J3) node[midway, right]{\small $a_3$};
                
                \filldraw[black] (O) circle (0.5pt);
                \filldraw[black] (J1) circle (0.5pt);
                \filldraw[black] (J2) circle (0.5pt);
                \filldraw[red] (J3) circle (0.5pt) node[above]{\small Brida};
            \end{tikzpicture}
            \end{center}
            \textbf{Objetivo:} Determine de forma independiente la pose del efector final respecto al marco base.
        \end{column}
    \end{columns}
\end{frame}

\begin{frame}{Ruta de Diagnóstico de Errores}
    \textit{Si su resultado fue incorrecto, ubique el punto de falla en la siguiente tubería[cite: 5]:}
    
    \vspace{0.5cm}
    \begin{center}
    \begin{tikzpicture}[
        node distance=0.2cm and 0.3cm,
        box/.style={rectangle, draw=ucbblue, thick, text width=1.5cm, align=center, rounded corners, minimum height=0.8cm, font=\scriptsize},
        arrow/.style={-Stealth, thick, ucbblue}
    ]
        \node[box] (a) {\textbf{A}\\Interpretación};
        \node[box, right=of a] (b) {\textbf{B}\\DH};
        \node[box, right=of b] (c) {\textbf{C}\\Matriz $T_i$};
        \node[box, right=of c] (d) {\textbf{D}\\Cadena};
        \node[box, right=of d] (e) {\textbf{E}\\Álgebra};
        \node[box, right=of e] (f) {\textbf{F}\\Pose};
        \node[box, right=of f] (g) {\textbf{G}\\Verificación};

        \draw[arrow] (a) -- (b);
        \draw[arrow] (b) -- (c);
        \draw[arrow] (c) -- (d);
        \draw[arrow] (d) -- (e);
        \draw[arrow] (e) -- (f);
        \draw[arrow] (f) -- (g);
    \end{tikzpicture}
    \end{center}
    
    \vspace{0.3cm}
    \begin{itemize}
        \item \textbf{[A-B] Concepto:} ¿Identificó bien $a_i, d_i, \alpha_i, \theta_i$?
        \item \textbf{[C-D] Estructura:} ¿Pudo armar ${}^{i-1}T_i$ y ordenarlas como $0 \to 1 \to 2 \to 3$?[cite: 5]
        \item \textbf{[E] Ejecución:} Errores aritméticos aislados (no críticos para el concepto).
        \item \textbf{[F-G] Físico:} ¿Es capaz de extraer $R$ y $p$ y comprobar si tienen sentido?[cite: 5]
    \end{itemize}
\end{frame}

\begin{frame}{Recordatorio Técnico: Matriz D-H Estándar}
    La convención del curso para una fila D-H es[cite: 5, 6]:
    \begin{equation*}
        {}^{i-1}T_i = R_z(\theta_i) T_z(d_i) T_x(a_i) R_x(\alpha_i)
    \end{equation*}

    Su expansión analítica directa es[cite: 5]:
    \begin{equation*}
        \boxed{
        {}^{i-1}T_i =
        \begin{bmatrix}
        c_{\theta_i} & -s_{\theta_i}c_{\alpha_i} & s_{\theta_i}s_{\alpha_i} & a_i c_{\theta_i} \\
        s_{\theta_i} & c_{\theta_i}c_{\alpha_i} & -c_{\theta_i}s_{\alpha_i} & a_i s_{\theta_i} \\
        0 & s_{\alpha_i} & c_{\alpha_i} & d_i \\
        0 & 0 & 0 & 1
        \end{bmatrix}
        }
    \end{equation*}
    
    \begin{alertblock}{Intervención Rápida}
        Si $(\theta, d, a, \alpha) = (q, 0.20, 0, -\pi/2)$, ¿cuál es el giro en el eslabón y cuál la traslación a lo largo de $Z$?[cite: 5]
    \end{alertblock}
\end{frame}

\begin{frame}{Semántica de Marcos y Extracción de Pose}
    \textbf{1. Encadenamiento Correcto}[cite: 5] \\
    El producto de transformaciones sigue una cadena semántica estricta:
    \begin{center}
        $\{0\} \xrightarrow{{}^0T_1} \{1\} \xrightarrow{{}^1T_2} \{2\} \xrightarrow{{}^2T_3} \{3\} \quad \implies \quad {}^0T_3 = {}^0T_1 {}^1T_2 {}^2T_3$
    \end{center}
    
    \vspace{0.4cm}
    \textbf{2. Anatomía de la Pose Final}[cite: 5] \\
    No basta con calcular la matriz; debe saber leerla:
    \begin{equation*}
        T = 
        \begin{bmatrix}
        \color{red}\mathbf{R} & \color{blue}\mathbf{p} \\
        \mathbf{0} & 1
        \end{bmatrix}
        \qquad \text{Ejemplo:} \quad
        {}^0T_2 = \begin{bmatrix} 0&0&1&\color{blue}0.20 \\ 1&0&0&\color{blue}-0.10 \\ 0&1&0&\color{blue}0.30 \\ 0&0&0&1 \end{bmatrix} \implies {}^0p_2 = \begin{bmatrix} \color{blue}0.20 \\ \color{blue}-0.10 \\ \color{blue}0.30 \end{bmatrix}
    \end{equation*}
\end{frame}

\begin{frame}{Problema Representativo 2: Transferencia (Espacial)}
    Considere un manipulador serial 3R con la siguiente tabla DH[cite: 5]:
    
    \begin{columns}
        \begin{column}{0.4\textwidth}
            \resizebox{\textwidth}{!}{
            \begin{tabular}{c c c c c}
            \toprule
            $i$ & $\theta_i$ & $d_i$ [m] & $a_i$ [m] & $\alpha_i$ \\ \midrule
            1 & $q_1$ & 0.20 & 0 & $+\pi/2$ \\
            2 & $q_2$ & 0 & 0.30 & 0 \\
            3 & $q_3$ & 0 & 0.20 & 0 \\ \bottomrule
            \end{tabular}
            }
            \vspace{0.3cm}
            
            \textbf{Variables:} \\
            $q_1=0^\circ, \ q_2=90^\circ, \ q_3=-90^\circ$[cite: 5].
        \end{column}
        
        \begin{column}{0.6\textwidth}
            \textbf{Tareas explícitas:}[cite: 5]
            \begin{enumerate}
                \item Construya las 3 transformaciones.
                \item Escriba la cadena de composición.
                \item Calcule ${}^0T_3$.
                \item Extraiga posición y orientación.
                \item Aplique una verificación física/matemática.
            \end{enumerate}
            \textit{Nota: Este problema introduce un giro en el primer eslabón, volviendo el sistema espacial[cite: 5].}
        \end{column}
    \end{columns}
\end{frame}

\begin{frame}{Lista de Verificación (Checklist) de la Pose Final}
    \textit{Antes de entregar su examen, aplique siempre estos filtros mentales a su matriz final[cite: 5]:}
    
    \vspace{0.4cm}
    \begin{itemize}
        \item[\checkmark] \textbf{Fila Homogénea:} ¿La última fila es estrictamente $[0 \ 0 \ 0 \ 1]$?[cite: 5]
        \item[\checkmark] \textbf{Rotación Ortonormal:} ¿La submatriz $R$ cumple $R^TR = I$? ¿Sus columnas son perpendiculares entre sí?[cite: 5]
        \item[\checkmark] \textbf{Límites Geométricos:} ¿La magnitud de la posición calculada $\|p\|$ es menor o igual a la suma de todos los eslabones ($a_i + d_i$)?[cite: 5]
        \item[\checkmark] \textbf{Dimensionalidad:} Si el robot es estrictamente planar en $XY$, ¿el componente $p_z$ dio exactamente $0$?[cite: 5]
    \end{itemize}
\end{frame}

\end{document}
